% This is text associated with the open textbook "Open Electromagnetics"
% See immediately below for author, date
% This work is licensed under the Creative Commons Attribution-ShareAlike 4.0 International License. (CC BY-SA 4.0) 
% To view a copy of the license, visit http://creativecommons.org/licenses/by-sa/4.0/

%ORB TITLE Magnetic Vector Potential
%ORB AUTHOR 0        # S.W. Ellingson, Virginia Tech (ellingson@vt.edu)
%ORB DATE 20191123   # yyyymmdd
%ORB ABSTRACT 

\label{m0195_Magnetic_Vector_Potential} % <-- Must match name of this file and module directory name.  Don't mess with this!
\index{magnetic vector potential|(}

A common problem in electromagnetics is to determine the fields radiated by a specified current distribution.
This problem can be solved using Maxwell's equations
\index{Maxwell's equations} 
along with the appropriate electromagnetic boundary conditions.
For time-harmonic 
\index{time-harmonic}
(sinusoidally-varying) currents, we use phasor representation.\footnote{Recall that there is no loss of generality in doing so, since any other time-domain variation in the current distribution can be represented using sums of time-harmonic solutions via the Fourier transform.}
Given the specified current distribution $\widetilde{\bf J}$ and the desired electromagnetic fields $\widetilde{\bf E}$ and $\widetilde{\bf H}$, the
appropriate equations are:
%
\begin{equation}
\nabla \cdot \widetilde{\bf E} = \widetilde{\rho}_v/\epsilon
\label{m0195_eMDE0}
\end{equation}
%
\begin{equation}
\nabla \times  \widetilde{\bf E} = -j\omega\mu\widetilde{\bf H} 
\label{m0195_eMCE}
\end{equation}
%
\begin{equation}
\nabla \cdot \widetilde{\bf H} = 0 
\label{m0195_eMDH}
\end{equation}
%
\begin{equation}
\nabla \times \widetilde{\bf H}  = \widetilde{\bf J} + j\omega\epsilon\widetilde{\bf E}
\label{m0195_eMCH} 
\end{equation}
%
where $\widetilde{\rho}_v$ is the volume charge density.
In most engineering problems, one is concerned with propagation through media which are well-modeled as homogeneous media with neutral charge, such as free space.\footnote{A counter-example would be propagation through a plasma,
\index{plasma}
which by definition consists of non-zero net charge.}
Therefore, in this section, we shall limit our scope to problems in which $\widetilde{\rho}_v=0$. 
Thus, Equation~\ref{m0195_eMDE0} simplifies to:
%
\begin{equation}
\nabla \cdot \widetilde{\bf E} = 0
\label{m0195_eMDE}
\end{equation}
%
To solve the linear system of partial differential Equations \ref{m0195_eMCE}--\ref{m0195_eMDE}, it is useful to invoke the concept of \emph{magnetic vector potential}.
The magnetic vector potential is a vector field that has the useful property that it is able to represent both the electric and magnetic fields as a single field.
This allows the formidable system of equations identified above to be reduced to a single equation which is simpler to solve.
Furthermore, this single equation turns out to be the wave equation, 
\index{wave equation!electromagnetic}
with the slight difference that the equation will be mathematically inhomogeneous, with the inhomogeneous part representing the source current.

The magnetic vector potential $\widetilde{\bf A}$ is defined by the following relationship:
%
\begin{equation}
\boxed{
\widetilde{\bf B} \triangleq \nabla \times \widetilde{\bf A}
}
\label{m0195_eMVPdef}
\end{equation}
%
where $\widetilde{\bf B}=\mu\widetilde{\bf H}$ is the magnetic flux density.
\index{magnetic flux density}
The magnetic field appears in three of Maxwell's equations. 
For Equation~\ref{m0195_eMVPdef} to be a reasonable definition, $\nabla \times \widetilde{\bf A}$ must yield reasonable results when substituted for $\mu\widetilde{\bf H}$ in each of these equations.
Let us first check for consistency with Gauss' law for magnetic fields,
\index{Gauss' law!magnetic field}
Equation~\ref{m0195_eMDH}.
Making the substitution, we obtain:
%
\begin{equation}
\nabla \cdot  \left( \nabla \times \widetilde{\bf A} \right) = 0
\end{equation}
%
This turns out to be a mathematical identity that applies to \emph{any} vector field (see Equation~\ref{m0140_eNdNxA} in Appendix~\ref{m0140_Vector_Identities}).
Therefore, Equation~\ref{m0195_eMVPdef} is consistent with Gauss' law for magnetic fields.

Next we check for consistency with Equation~\ref{m0195_eMCE}.  Making the substitution:
%
\begin{equation}
\nabla \times  \widetilde{\bf E} = -j\omega\left(\nabla \times \widetilde{\bf A}\right)
\end{equation}
%
Gathering terms on the left, we obtain
%
\begin{equation}
\nabla \times  \left( \widetilde{\bf E} +j\omega \widetilde{\bf A} \right) = 0
\label{m0195_ephi1}
\end{equation}
%
Now, for reasons that will become apparent in just a moment, we define a new scalar field $\widetilde{V}$ and require it to satisfy the following relationship:
%
\begin{equation}
-\nabla\widetilde{V} \triangleq \widetilde{\bf E} +j\omega \widetilde{\bf A}
\label{m0195_ephi2} 
\end{equation}
%
Using this definition, Equation~\ref{m0195_ephi1} becomes:
%
\begin{equation}
\nabla \times  \left( -\nabla\widetilde{V} \right) = 0
\end{equation}
%
which is simply
%
\begin{equation}
\nabla \times  \nabla\widetilde{V} = 0
\end{equation}
%
Once again we have obtained a mathematical identity that applies to \emph{any} vector field (see Equation~\ref{m0140_eNxNA} in Appendix~\ref{m0140_Vector_Identities}).
Therefore, $\widetilde{V}$ can be \emph{any} mathematically-valid scalar field.
Subsequently, Equation~\ref{m0195_eMVPdef} is consistent with Equation~\ref{m0195_eMCE} (Maxwell's curl equation for the electric field) for any choice of $\widetilde{V}$ that we are inclined to make. 

Astute readers might already realize what we're up to here.  
Equation~\ref{m0195_ephi2} is very similar to the relationship ${\bf E} = -\nabla V$ from electrostatics,\footnote{Note: No tilde in this expression.} in which $V$ is the scalar electric potential field.
Evidently Equation~\ref{m0195_ephi2} is an enhanced version of that relationship that accounts for the coupling with ${\bf H}$ (here, represented by ${\bf A}$) in the time-varying (decidedly non-static) case.
That assessment is correct, but let's not get too far ahead of ourselves:  As demonstrated in the previous paragraph, we are not yet compelled to make any particular choice for $\widetilde{V}$, and this freedom will be exploited later in this section.

Next we check for consistency with Equation~\ref{m0195_eMCH}.  Making the substitution:
%
\begin{equation}
\nabla \times \left(\frac{1}{\mu}\nabla \times \widetilde{\bf A}\right) = \widetilde{\bf J} +j\omega\epsilon\widetilde{\bf E}
\end{equation}
%
Multiplying both sides of the equation by $\mu$:
%
\begin{equation}
\nabla \times \nabla \times \widetilde{\bf A} = \mu\widetilde{\bf J} +j\omega\mu\epsilon\widetilde{\bf E}
\end{equation}
%
Next we use Equation~\ref{m0195_ephi2} to eliminate $\widetilde{\bf E}$, yielding:
%
\begin{equation}
\nabla \times \nabla \times \widetilde{\bf A} = \mu\widetilde{\bf J} +j\omega\mu\epsilon\left(-\nabla\widetilde V -j\omega \widetilde{\bf A}\right)
\end{equation}
%
After a bit of algebra, we obtain
%
\begin{equation}
\nabla \times \nabla \times \widetilde{\bf A}
= \omega^2\mu\epsilon\widetilde{\bf A} - j\omega\mu\epsilon\nabla\widetilde V + \mu\widetilde{\bf J} 
\label{m0195_e1}
\end{equation}
%
Now we replace the left side of this equation using vector identity \ref{m0140_eNxNxA} in Appendix~\ref{m0140_Vector_Identities}:
%
\begin{equation}
\nabla \times \nabla \times \widetilde{\bf A} \equiv \nabla\left(\nabla \cdot \widetilde{\bf A}\right) - \nabla^2 \widetilde{\bf A}
\end{equation}
%
Equation~\ref{m0195_e1} becomes:
%
\begin{equation}
\nabla\left(\nabla \cdot \widetilde{\bf A}\right) - \nabla^2 \widetilde{\bf A}
= \omega^2\mu\epsilon\widetilde{\bf A} - j\omega\mu\epsilon\nabla\widetilde V + \mu\widetilde{\bf J} 
\end{equation}
%
Now multiplying both sides by $-1$ and rearranging terms:
%
\begin{equation}
\nabla^2 \widetilde{\bf A}
+\omega^2\mu\epsilon\widetilde{\bf A} 
= 
\nabla\left(\nabla \cdot \widetilde{\bf A}\right)
+j\omega\mu\epsilon\nabla\widetilde V 
-\mu\widetilde{\bf J} 
\end{equation}
%
Combining terms on the right side:
%
\begin{equation}
\nabla^2 \widetilde{\bf A}
+\omega^2\mu\epsilon\widetilde{\bf A} 
= 
\nabla\left( \nabla \cdot \widetilde{\bf A} +j\omega\mu\epsilon\widetilde V\right) 
-\mu\widetilde{\bf J} 
\label{m0195_eWEA1}
\end{equation}
%
Now consider the expression
$\nabla \cdot \widetilde{\bf A} +j\omega\mu\epsilon\widetilde V$
appearing in the parentheses on the right side of the equation.
We established earlier that $\widetilde V$ can be essentially any scalar field -- from a mathematical perspective, we are free to choose.
Invoking this freedom, we now require $\widetilde V$ to satisfy the following expression:
%
\begin{equation}
\nabla \cdot \widetilde{\bf A} +j\omega\mu\epsilon\widetilde V = 0
\label{m0195_eLGC}
\end{equation}
%
Clearly this is advantageous in the sense that Equation~\ref{m0195_eWEA1} is now dramatically simplified. This equation becomes:
%
\begin{equation}
\boxed{
\nabla^2 \widetilde{\bf A}
+\omega^2\mu\epsilon\widetilde{\bf A} 
= 
-\mu\widetilde{\bf J} 
}
\label{m0195_ePDEA}
\end{equation}
%
Note that this expression is a \emph{wave equation}.
\index{wave equation!electromagnetic}
\index{wave equation!magnetic vector potential}
In fact it is the \emph{same} wave equation that determines $\widetilde{\bf E}$ and $\widetilde{\bf H}$ in source-free regions, \emph{except} the right-hand side is not zero.
Using mathematical terminology, we have obtained an equation for $\widetilde{\bf A}$ in the form of an \emph{inhomogeneous} partial differential equation, where the inhomogeneous part includes -- no surprise here -- the source current $\widetilde{\bf J}$.

Now we have what we need to find the electromagnetic fields radiated by a current distribution.
The procedure is simply as follows:
%
\begin{enumerate}
\item Solve the partial differential Equation~\ref{m0195_ePDEA} for $\widetilde{\bf A}$ along with the appropriate electromagnetic boundary conditions.
\item $\widetilde{\bf H} = (1/\mu) \nabla \times \widetilde{\bf A}$
\item $\widetilde{\bf E}$ may now be determined from $\widetilde{\bf H}$ using Equation~\ref{m0195_eMCH}.
\end{enumerate}

Summarizing:
%
%%%%%%%%%%%%%%%%%%%%%%%%%%%%%%%%%%%%%%%%%%%%%%%
%ORB HIGHLIGHT
\begin{mdframed}[backgroundcolor=yellow!20,nobreak=true] 
The magnetic vector potential $\widetilde{\bf A}$ is a vector field, defined by Equation~\ref{m0195_eMVPdef}, that is able to represent both the electric and magnetic fields simultaneously.
\end{mdframed}
%%%%%%%%%%%%%%%%%%%%%%%%%%%%%%%%%%%%%%%%%%%%%%%
%
Also:
%
%%%%%%%%%%%%%%%%%%%%%%%%%%%%%%%%%%%%%%%%%%%%%%%
%ORB HIGHLIGHT
\begin{mdframed}[backgroundcolor=yellow!20,nobreak=true] 
To determine the electromagnetic fields radiated by a current distribution $\widetilde{\bf J}$, one may solve Equation~\ref{m0195_ePDEA} for $\widetilde{\bf A}$ and then use Equation~\ref{m0195_eMVPdef} to determine $\widetilde{\bf H}$ and subsequently $\widetilde{\bf E}$.
\end{mdframed}
%%%%%%%%%%%%%%%%%%%%%%%%%%%%%%%%%%%%%%%%%%%%%%%
%
Specific techniques for performing this procedure -- in particular, for solving the differential equation -- vary depending on the problem, and are discussed in other sections of this book.

We conclude this section with a few comments about Equation~\ref{m0195_eLGC}.
This equation is known as the \emph{Lorenz gauge condition}.
\index{Lorenz gauge condition}
This constraint is not quite as arbitrary as the preceding derivation implies; rather, there is some deep physics at work here.
Specifically, the Lorenz gauge leads to the classical interpretation of $\widetilde{V}$ as the familiar scalar electric potential, as noted previously in this section. (For additional information on that idea, recommended starting points are included in ``Additional Reading'' at the end of this section.)

At this point, it should be clear that the electric and magnetic fields are not merely coupled quantities, but in fact two aspects of the same field; namely, the magnetic vector potential.
In fact modern physics (quantum mechanics) yields the magnetic vector potential as a description of the ``electromagnetic force,'' a single entity which constitutes one of the four fundamental forces recognized in modern physics; the others being gravity, the strong nuclear force, and the weak nuclear force. 
For more information on that concept, an excellent starting point is the video ``Quantum Invariance \& The Origin of The Standard Model'' referenced at the end of this section.

%%%%%%%%%%%%%%%%%%%%%

{\bf Additional Reading:}
\begin{itemize}
\item ``\href{https://en.wikipedia.org/wiki/Lorenz_gauge_condition}{Lorenz gauge condition}'' on Wikipedia.
\item ``\href{https://en.wikipedia.org/wiki/Magnetic_potential}{Magnetic potential}'' on Wikipedia.
\item PBS Space Time video ``\href{https://youtu.be/V5kgruUjVBs}{Quantum Invariance} \& The Origin of The Standard Model,'' available on YouTube.
\end{itemize}

\index{magnetic vector potential|)}


